\subsection{Reading and Synthesizing Research Papers}

Reading depth should match the decision at hand: triage, comprehension, verification, reproduction, or synthesis.

\begin{process} \label{proc:three-pass-paper-reading} \hlt{Three-Pass Paper Reading}\\
Adapt the three-pass method of \citet{keshav2007readpaper} to the paper's evidence type.
\begin{enumerate}[label=(\roman*)]
\setlength{\itemsep}{0pt}
\item \hlt{Triage}: inspect title, abstract, introduction, headings, conclusion, figures, and references. Record category, context, claimed contribution, apparent assumptions, and relevance.
\item \hlt{Comprehension}: read the body carefully, map each main claim to its stated evidence, inspect figures and tables, and mark unfamiliar prerequisites or missing details.
\item \hlt{Reconstruction}: independently rebuild the proof, derivation, algorithm, or experiment closely enough to expose hidden assumptions, fragile steps, and reproducibility barriers.
\end{enumerate}
Stop after the earliest pass sufficient for the decision, and resume only after acquiring a missing prerequisite when further reading would otherwise be guesswork.
\end{process}

\begin{process} \label{proc:paper-claim-evidence-ledger} \hlt{Paper Claim--Evidence Ledger}\\
Create one entry per central claim and record its locator, exact scope, assumptions, evidence type, dependency, result, uncertainty, and limitation. For a mathematical paper, reconstruct definitions and proof obligations. For an empirical paper, reconstruct the sample, protocol, comparator, metric, variability, and analysis. For a systems paper, reconstruct the boundary, workload, implementation, platform, and failure model. Mark each claim as supported, partially supported, unsupported, or not assessed; reputation and venue are metadata rather than evidence.
\end{process}

\begin{process} \label{proc:literature-synthesis} \hlt{Literature Search and Synthesis}\\
Build a literature map rather than a pile of summaries.
\begin{enumerate}[label=(\roman*)]
\setlength{\itemsep}{0pt}
\item \hlt{Seed}: start from precise keywords, a recent survey, and several directly relevant papers.
\item \hlt{Backward search}: follow repeated foundational citations to original results, data, or methods.
\item \hlt{Forward search}: locate corrections, replications, extensions, and critiques that cite each key work.
\item \hlt{Cluster}: group papers by question, assumptions, method, evidence, and operating context rather than by author.
\item \hlt{Compare}: normalize terminology and record agreements, contradictions, gaps, and incomparable protocols.
\item \hlt{Saturate}: stop when new searches cease changing the conceptual map, while retaining the search date and boundary.
\end{enumerate}
Keshav's survey procedure likewise expands seed papers through repeated authors, shared references, and recent proceedings \citep{keshav2007readpaper}.
\end{process}

\begin{warning} \label{warn:paper-version-and-source-status} \hlt{Version and Source Status}\\
Before relying on a paper, check its version, venue, peer-review status, corrections, retractions, supplements, artifact links, and conflicts of interest. Cite the original source for a result or dataset when available. A preprint may be valuable, and peer review may improve reliability, but neither status settles the truth of a claim.
\end{warning}

\begin{principle} \label{princ:read-against-the-claim} \hlt{Read Against the Claim}\\
Evaluate the strongest wording in the title, abstract, or conclusion against the narrowest evidence in the body. The decisive questions are what was established, under which assumptions, by which evidence, with what uncertainty, and within which boundary.
\end{principle}
